\section{Introduction}

In this article, we study the nonlinear Kirchhoff-Schr\"odinger-Poisson system
with Sobolev critical growth
\begin{equation}\label{lab1.1}
\begin{gathered}
(a+b\int_{\mathbb{R}^3}|(-\Delta)^{s/2}u|^2dx)(-\Delta)^s u+ \phi u
= \lambda u+\mu|u|^{p-2}u+|u|^{{2^*_s}-2}u,  \quad \text{in }\mathbb{R}^3, \\
(-\Delta)^s \phi=u^2, \quad   \text{in }\mathbb{R}^3,
\end{gathered}
\end{equation}
where $a,b>0$, $s\in(\frac{3}{4},1)$, and  $p \in(2,2^*_s)$,  $\mu>0$ and
$\lambda \in \mathbb{R}$ are parameters.
Here $(-\Delta)^s$ $(s\in(0,1))$ is the fractional Laplacian operator which is
defined by
$$
(-\Delta)^s u(x)=C_s {\rm P.V.}\int_{\mathbb{R}^3}\frac{u(x)-u(y)}{|x-y|^{3+2s}}dy
=C_s\lim_{\epsilon\to 0}\int_{{\mathbb{R}^3}
\setminus B_\epsilon(x)}\frac{u(x)-u(y)}{|x-y|^{3+2s}}dy,
$$
for $u\in \mathcal{S}(\mathbb{R}^3)$, where $\mathcal{S}(\mathbb{R}^3) $ is the
Schwartz space  of rapidly decaying $C^\infty$ functions, $B_\epsilon(x)$ denote an
open ball of radius $\epsilon$ at $x$ and $C_s$ is a normalization constant.

System \eqref{lab1.1} has been motivated by the
time-dependent fractional Schr\"odinger-Poisson system
\begin{equation}\label{lab1.12}
\begin{gathered}
i\frac{\partial \Psi}{\partial \tau } =(-\Delta)^s\Psi +\lambda\phi\Psi-f(x,|\Psi|),
\quad  x\in \mathbb{R}^3,\\
(-\Delta)^t\phi=|\Psi|^2, \quad  x\in \mathbb{R}^3,
\end{gathered}
\end{equation}
where $\Psi: {\mathbb{R}\times \mathbb{R}^3}\to \mathbb{C}$, $s,t \in(0,1)$,
$\lambda \in \mathbb{R}$. It is well-known that, the first equation in
system \eqref{lab1.12} was used by Laskin (see \cite{L2998, L3998}) to extend the Feynman path integral, from Brownian-like to L$\mathrm {\acute{e}}$vy-like quantum mechanical paths. This class of fractional Schr\"odinger equations with a repulsive nonlocal Coulombic potential can be approximated by the Hartree-Fock equations to describe a quantum mechanical system of many particles. For more application backgrounds on the fractional Laplacian see \cite{C4998, D7998, M6998}.

When looking for solutions to system (1.1), there are two distinct options  regarding  the frequency parameter  $\lambda$. One is to regard the frequency $\lambda$ as a given constant. Xiang and Wang \cite{H2998} first investigated  the  fractional  Kirchhoff-Schr\"odinger-Poisson system, and obtained the existence, multiplicity and asymptotic  behavior  of  nonnegative  solutions. In recent years, researchers have shown growing interest in the fractional Kirchhoff-Schr\"odinger-Poisson system
\begin{equation}\label{lab1.156}
\begin{gathered}
(a+b\int_{\mathbb{R}^3}|(-\Delta)^{s/2}u|^2dx)(-\Delta)^s u+V(x)u+\mu \phi u= f(x,u),   
\quad  \text{in }\mathbb{R}^3, \\
(-\Delta)^t \phi=\mu u^2, \quad   \text{in }\mathbb{R}^3,
\end{gathered}
\end{equation}
where $a>0$, $b\geq 0$. When $V(x)=0$, $f(x,u)=f(u)$, by utilizing minimax argument, 
Ambrosio \cite{A1111} obtained the existence of a nontrivial solutions  
for system \eqref{lab1.156} with Berestycki-Lions type nonlinearities. 
When   $V(x)\neq 0$, Wang et.al \cite{W1112} studied the existence of ground solutions 
for system \eqref{lab1.156} with  $V(x)=1$ and 
$f(x,u)=(|x|^{-\theta}\ast F(u))f(u)$, $\theta \in(0, 3-2t)$, and used the Pohoz\v{a}ev 
type  manifold; Then, under some assumptions on $V$ and $f$, by using constraint 
variational approach and a quantitative deformation lemma, Meng et.al \cite{M1113} 
proved the existence of  the least energy sign-changing solutions for system 
\eqref{lab1.156}.  Feng et al.\ \cite{F1114} applied a similar method studied the 
least energy sign-changing solutions to the following critical fractional 
Kirchhoff-Schr\"odinger-Poisson system with steep potential well
  %\label{lab1.1156}
\begin{gather*}
\Big(a+b\int_{\mathbb{R}^3}|(-\Delta)^{s/2} u|^2dx\Big)
(-\Delta)^s u+V_{\lambda}(x)u+ \phi u=|u|^{p-2}u+|u|^{2^*_s-2}u, \quad
 \text{in }\mathbb{R}^3, \\
(-\Delta)^t \phi=u^2, \quad   \text{in }\mathbb{R}^3,
\end{gather*}
where $s\in(\frac{3}{4}, 1)$, $t\in(0,1)$, $V_\lambda(x)=\lambda V(x)+1$,
$\lambda>0$ and $p>4$.  Jian et al.\ \cite{J1115}, deal with the
fractional Kirchhoff-Schr\"odinger-Poisson system with steep potential well
 % \label{lab1.1567}
\begin{gather*}
(a+b\int_{\mathbb{R}^3}|(-\Delta)^{s/2}u|^2dx)(-\Delta)^s u+\lambda V(x)u
 +\mu \phi u= |u|^{p-2}u,  \quad \text{in }\mathbb{R}^3, \\
(-\Delta)^t \phi= u^2, \quad   \text{in }\mathbb{R}^3,
\end{gather*}
where $s\in[\frac{3}{4}, 1)$, $t\in(0,1)$, $2<p<4$, $a>0$ is a constant, and
$b, \lambda, \mu$ are positive parameters. By applying the truncation technique and
the parameter-dependent compactness lemma, they first proved the existence of
positive solutions.  Furthermore, they investigated the   asymptotic behavior as
$b\to 0$, $\lambda \to \infty$ and $\mu \to 0$, respectively.
For other existence results, we refer to \cite{F1125, W1145, W1135}  and the
references therein.

When $a=1$, $b=0$, system \eqref{lab1.156} reduces to the following fractional
Schr\"odinger-Poisson system
\begin{equation}\label{lab1.15776}
\begin{gathered}
(-\Delta)^s u+V(x)u+\mu \phi u= f(x,u),   \quad \text{in }\mathbb{R}^3, \\
(-\Delta)^t \phi=\mu u^2, \quad   \text{in }\mathbb{R}^3.
\end{gathered}
\end{equation}
Recently, under various potentials and nonlinear terms, most scholars have
investigated the existence and multiplicity of ground state  solutions,
sign-changing solutions, and nontrivial solutions for system \eqref{lab1.15776}.
 For further details, we refer  the interesting readers to see
\cite{L1159, D1399, G1312, G1314, P1199} and so on.

Alternatively, the other one is to regard the frequency $\lambda$ as an unknown
quantity. In such point of view, it is natural to prescribe the mass, i.e.,
the $L^2$-norm, so that $\lambda$ can be interpreted as a Lagrange multiplier.
Solutions of this type are often referred to as normalized solutions.
Nowadays, from a physical point of view, some physicists are very interested in the
 normalized solutions, see for example \cite{Bartsch99, luo288,  yang1288}.
There a few results are related to the study of normalized solutions for the
fractional Kirchhoff-Schr\"odinger-Poisson system except Wang et al.\ \cite{wang1388}.
They considered the  fractional Kirchhoff-Schr\"odinger-Poisson system
\[  %\label{lab1.1234}
\begin{gathered}
\Big(a+b\int_{\mathbb{R}^3}|(-\Delta)^{s/2}u|^2dx\Big)(-\Delta)^s u+ \phi u= f
(u)+\lambda u    \quad \text{in }\mathbb{R}^3, \\
(-\Delta)^t \phi=u^2 \quad    \text{in }\mathbb{R}^3,
\end{gathered}
\]
where $a,b>0$, $s, t\in (0,1)$, $2s+2t>3$, $\lambda\in \mathbb{R}$,
$f\in C(\mathbb{R}, \mathbb{R})$ satisfies some general conditions which contain
the case  $f(u)\sim|u|^{p-2}u$ with
$p\in(\frac{8s+4t-3}{s+t},\frac{8s}{3}+2)\cup(\frac{8s}{3}+2, 2^*_s)$,
$2^*_s=\frac{6}{3-2s}$. They obtained
the existence of normalized solutions by using the Pohoz\v{a}ev manifold and
variational method.

 Recently, He and Meng \cite{H2588} studied the existence and multiplicity of the
normalized solutions for the nonlinear fractional Schr\"odinger-Poisson system with
Sobolev critical exponent
\[ % \label{lab1.1434}
\begin{gathered}
(-\Delta)^s u+ \alpha\phi u= \mu|u|^{p-2}u+|u|^{2^*_s-2}u+\lambda u   \quad
 \text{in }\mathbb{R}^3, \\
(-\Delta)^t \phi=u^2 \quad    \text{in }\mathbb{R}^3,
\end{gathered}
\]
 where $s, t\in (0,1)$, $2s+2t>3$, $p\in(2, 2^*_s)$, $\alpha, \mu>0$ are parameters
and $\lambda\in \mathbb{R}$.

Motivated by \cite{H2588} and \cite{wang1388}, a natural question is whether the
fractional Kirchhoff-Schr\"odinger-Poisson system with the Sobolev critical growth can
be applied to obtain the  existence and multiplicity   of normalized solutions for $p$
in distinct ranges.
Therefore, we study system \eqref{lab1.1} and give an affirmative answer.
In addition, to recover the compactness, we will take $H^s_r(\mathbb{R}^3)$ as a
working space.

By using the Lax-Milgram theorem, for any $u\in H^s_r(\mathbb{R}^3)$, a unique
$\phi ^s_u(x)\in D^{s,2}(\mathbb{R}^3)$ is given by
\[ \label{lab1.4}
\phi ^s_u(x)=|x|^{2s-3}\ast|u|^2=\int_{\mathbb{R}^3}\frac{u^2(y)}{|x-y|^{3-2s}}dy,
\]
such that $(-\Delta)^s\phi =u^2$ and that inserting it into the first equation of
system \eqref{lab1.1}, then system \eqref{lab1.1} can be transformed into the following
single equation
\begin{equation}\label{lab1.5}
\Big(a+b\int_{\mathbb{R}^3}|(-\Delta)^{s/2}u|^2dx\Big)(-\Delta)^s u
 + \phi^s_u u= \lambda u+ \mu|u|^{p-2}u+ |u|^{2^*_s-2}u, \quad\text{in }\mathbb{R}^3.
\end{equation}
It can be proved  that to find the normalized solutions of system \eqref{lab1.1}
is to seek the critical points of the functional
\begin{align*}
  I_\mu(u)
&=\frac{a}{2}\int_{\mathbb{R}^3}|(-\Delta)^{s/2}u|^2dx
  +\frac{b}{4}\Big(\int_{\mathbb{R}^3}|(-\Delta)^{s/2}u|^2dx\Big)^2\\
&\quad +\frac{1}{4}\int_{\mathbb{R}^3}\phi ^s_uu^2dx
 -\frac{\mu}{p}\int _{\mathbb{R}^3}|u|^pdx
 -\frac{1}{2^*_s}\int _{\mathbb{R}^3}|u|^{2^*_s}dx,
\end{align*}
under the constraint
\[ \label{1.8}
S_r(c)=\Big\{u \in H^s_r(\mathbb{R}^3):\int_{\mathbb{R}^3}|u|^2dx=c^2,c>0\Big\}.
\]
It is well known that $I_\mu\in {C}^1(H^s_r(\mathbb{R}^3), \mathbb{R})$.
Here are our main results.

